%------------------------------------------------------------------------------%
\addtocounter{exemplo}{1}
\begin{frame}{Exemplo \theexemplo}
  Encontre os pontos críticos e valores extremos locais da função
  \[
     f(x, y) = x^2 + y^2 - 4y + 9
  \]
\end{frame}

%------------------------------------------------------------------------------%
\begin{frame}{Exemplo \theexemplo}
  Derivadas parciais
  \pause
  \vspace{\baselineskip}
  \[
    \D{f}{x}
    \pause
    = \D{}{x}\left(x^2 + y^2 - 4y + 9\right)
    \pause
    = 2x
  \]
  \pause
  \vspace{\baselineskip}
  \[
    \D{f}{y}
    \pause
    = \D{}{y}\left(x^2 + y^2 - 4y + 9\right)
    \pause
    = 2y -4
  \]
\end{frame}

%------------------------------------------------------------------------------%
\begin{frame}{Exemplo \theexemplo}
  \onslide<+->{Pontos críticos}

  \vspace{0.1\baselineskip}
  \onslide<+->{As derivadas existem em todos os pontos do plano}

  \vspace{0.1\baselineskip}
  \onslide<+->{Plano tangente horizontal}
  \begin{columns}
  \column{0.33\linewidth}
    \begin{align*}
      \onslide<+->{f_x(x, y) & = 0 \\}
      \onslide<+->{2x        & = 0 \\}
      \onslide<+->{x         & = 0}
    \end{align*}
  \column{0.33\linewidth}
    \begin{align*}
      \onslide<+->{f_y(x, y) & = 0 \\}
      \onslide<+->{2y -4     & = 0 \\}
      \onslide<+->{y         & = 2}
    \end{align*}
    \column{0.33\linewidth}
    \onslide<+->{Ponto critico é \((0, 2)\)}
  \end{columns}
\end{frame}

%------------------------------------------------------------------------------%
\begin{frame}{Exemplo \theexemplo}
  Onde a função vale
  \[
    f(0, 2)
    \pause
    = \left(x^2 + y^2 - 4y + 9\right)\evalat{(0, 2)}{}
    \pause
    = 0 + 4 - 8 + 9
    \pause
    = 5
  \]

  \pause
  \vspace{\baselineskip}
  Analisando a função percebemos que 5 é o menor valor que ela assume,
  portanto é um ponto de mínimo local (e nesse caso global também)
\end{frame}

%------------------------------------------------------------------------------%
